% GENERATED FILE. Edit canonical lessons, then run npm run generate:books.
% canonical-source-hash: fnv1a64:d6c4368ab12ef5cc
% canonical-lessons: ML-C401-atima, ML-C401-viplavam, ML-C401-adhinivesam, ML-C401-poratuka, ML-C401-ceruttunilkkuka

\chapter{Invasion and Revolt}
\label{ch:ml-invasion-and-revolt}
\hlchaptermodality{401}

\begin{chapteropening}
\textbf{By the end of this chapter:} \emph{I can say a slave, a revolution, an invasion, to fight and to resist.}

Complete the last lesson of chapter 401: I can say a slave, a revolution, an invasion, to fight and to resist.
\end{chapteropening}

\section[\texorpdfstring{\ml{അടിമ} (aṭima)}{അടിമ (aṭima)}]{\ml{അടിമ} (aṭima) — a slave}

\label{lesson:ML-C401-atima}



\begin{quote}
Before the new one: say the Malayalam for to beat in battle, then the Malayalam for to surrender.
\end{quote}

\subsection*{You'll want to know: \ml{അടിമ}}
\textbf{\ml{അടിമ}} — \emph{aṭima} — \textquotedblleft{}a slave\textquotedblright{}.

\begin{quote}
\ml{പണ്ട്} \ml{ആളുകളെ} \ml{അടിമകളായി} \ml{വിറ്റിരുന്നു}.
\end{quote}

\begin{quote}
\emph{paṇṭŭ āḷukaḷe aṭimakaḷāyi viṟṟirunnu.}
\end{quote}

\begin{quote}
Long ago people were sold as slaves.
\end{quote}

Plural: \ml{അടിമകൾ}, \emph{aṭimakaḷ}.

Plural \ml{അടിമകൾ}; slavery is \ml{അടിമത്തം}.

\begin{cousinweb}
\textbf{\ml{അടിമ}} (\emph{aṭima}) is native Dravidian: DEDR 72 lists Malayalam \emph{aṭima}, \textquotedblleft{}slavery, slave\textquotedblright{}, under \ml{അടി}, \textquotedblleft{}the sole of the foot, the bottom\textquotedblright{}: one who is kept at the feet of a master.
\end{cousinweb}

\begin{grammarlens}[title={What you've built}]
One more word for telling what happened.
\end{grammarlens}

\subsection*{Your turn}
\begin{itemize}
\raggedright
  \item \emph{Say it:} \emph{aṭima}
  \item \emph{Say it:} \emph{aṭima}, once more
  \item \emph{Say it:} \emph{paṇṭŭ āḷukaḷe aṭimakaḷāyi viṟṟirunnu.} — Long ago people were sold as slaves.
  \item \emph{Recall:} say the Malayalam for to beat in battle, then the Malayalam for to surrender, then say \emph{aṭima} again
\end{itemize}

\begin{tcolorbox}[breakable,colback=teal!4,colframe=teal!35!black,title={Before you move on}]
What does \ml{അടിമ} mean? (\textquotedblleft{}A slave\textquotedblright{}.) Say it once more.
\end{tcolorbox}
\section[\texorpdfstring{\ml{വിപ്ലവം} (viplavaṁ)}{വിപ്ലവം (viplavaṁ)}]{\ml{വിപ്ലവം} (viplavaṁ) — a revolution}

\label{lesson:ML-C401-viplavam}



\begin{quote}
Before the new one: say the Malayalam for to surrender, then the Malayalam for a slave.
\end{quote}

\subsection*{You'll want to know: \ml{വിപ്ലവം}}
\textbf{\ml{വിപ്ലവം}} — \emph{viplavaṁ} — \textquotedblleft{}a revolution\textquotedblright{}.

\begin{quote}
\ml{ആ} \ml{വിപ്ലവം} \ml{രാജ്യത്തെ} \ml{മാറ്റി}.
\end{quote}

\begin{quote}
\emph{ā viplavaṁ rājyatte māṟṟi.}
\end{quote}

\begin{quote}
That revolution changed the country.
\end{quote}

Also any great upheaval in a field, as in farming or technology.

\begin{cousinweb}
\textbf{\ml{വിപ്ലവം}} (\emph{viplavaṁ}) is borrowed from Sanskrit \emph{viplava}, \textquotedblleft{}confusion, trouble, calamity\textquotedblright{} (Monier-Williams); modern Malayalam uses it for a revolution.
\end{cousinweb}

\begin{grammarlens}[title={What you've built}]
One more word for telling what happened.
\end{grammarlens}

\subsection*{Your turn}
\begin{itemize}
\raggedright
  \item \emph{Say it:} \emph{viplavaṁ}
  \item \emph{Say it:} \emph{viplavaṁ}, once more
  \item \emph{Say it:} \emph{ā viplavaṁ rājyatte māṟṟi.} — That revolution changed the country.
  \item \emph{Recall:} say the Malayalam for to surrender, then the Malayalam for a slave, then say \emph{viplavaṁ} again
\end{itemize}

\begin{tcolorbox}[breakable,colback=teal!4,colframe=teal!35!black,title={Before you move on}]
What does \ml{വിപ്ലവം} mean? (\textquotedblleft{}A revolution\textquotedblright{}.) Say it once more.
\end{tcolorbox}
\section[\texorpdfstring{\ml{അധിനിവേശം} (adhinivēśaṁ)}{അധിനിവേശം (adhinivēśaṁ)}]{\ml{അധിനിവേശം} (adhinivēśaṁ) — an invasion; occupation}

\label{lesson:ML-C401-adhinivesam}



\begin{quote}
Before the new one: say the Malayalam for a slave, then the Malayalam for a revolution.
\end{quote}

\subsection*{You'll want to know: \ml{അധിനിവേശം}}
\textbf{\ml{അധിനിവേശം}} — \emph{adhinivēśaṁ} — \textquotedblleft{}an invasion; occupation\textquotedblright{}.

\begin{quote}
\ml{വിദേശ} \ml{അധിനിവേശം} \ml{നൂറ്} \ml{വർഷം} \ml{നീണ്ടു}.
\end{quote}

\begin{quote}
\emph{vidēśa adhinivēśaṁ nūṟŭ varṣaṁ nīṇṭu.}
\end{quote}

\begin{quote}
The foreign occupation lasted a hundred years.
\end{quote}

A learned word of history books and the news.

\begin{cousinweb}
\textbf{\ml{അധിനിവേശം}} (\emph{adhinivēśaṁ}) is a learned word usually traced to Sanskrit \emph{adhi}, \textquotedblleft{}over, upon\textquotedblright{}, and \emph{niveśa}, \textquotedblleft{}settling, a camp\textquotedblright{}: to settle on top of someone else's land.
\end{cousinweb}

\begin{grammarlens}[title={What you've built}]
One more word for telling what happened.
\end{grammarlens}

\subsection*{Your turn}
\begin{itemize}
\raggedright
  \item \emph{Say it:} \emph{adhinivēśaṁ}
  \item \emph{Say it:} \emph{adhinivēśaṁ}, once more
  \item \emph{Say it:} \emph{vidēśa adhinivēśaṁ nūṟŭ varṣaṁ nīṇṭu.} — The foreign occupation lasted a hundred years.
  \item \emph{Recall:} say the Malayalam for a slave, then the Malayalam for a revolution, then say \emph{adhinivēśaṁ} again
\end{itemize}

\begin{tcolorbox}[breakable,colback=teal!4,colframe=teal!35!black,title={Before you move on}]
What does \ml{അധിനിവേശം} mean? (\textquotedblleft{}An invasion; occupation\textquotedblright{}.) Say it once more.
\end{tcolorbox}
\section[\texorpdfstring{\ml{പോരാടുക} (pōrāṭuka)}{പോരാടുക (pōrāṭuka)}]{\ml{പോരാടുക} (pōrāṭuka) — to fight; to struggle}

\label{lesson:ML-C401-poratuka}



\begin{quote}
Before the new one: say the Malayalam for a revolution, then the Malayalam for an invasion.
\end{quote}

\subsection*{You'll want to know: \ml{പോരാടുക}}
\textbf{\ml{പോരാടുക}} — \emph{pōrāṭuka} — \textquotedblleft{}to fight; to struggle\textquotedblright{}.

\begin{quote}
\ml{നാട്ടുകാർ} \ml{ശത്രുക്കളോട്} \ml{പോരാടി}.
\end{quote}

\begin{quote}
\emph{nāṭṭukār śatrukkaḷōṭŭ pōrāṭi.}
\end{quote}

\begin{quote}
The local people fought the enemies.
\end{quote}

Past: \ml{പോരാടി}, \emph{pōrāṭi}.

Also a struggle for a cause: \ml{അവകാശത്തിനുവേണ്ടി} \ml{പോരാടുക}.

\begin{cousinweb}
\textbf{\ml{പോരാടുക}} (\emph{pōrāṭuka}) is native Dravidian: DEDR 4540 lists Malayalam \emph{pōr}, \textquotedblleft{}battle, war\textquotedblright{}, and \emph{pōr-āṭuka}, \textquotedblleft{}to fight, to contend\textquotedblright{}.
\end{cousinweb}

\begin{grammarlens}[title={What you've built}]
One more word for telling what happened.
\end{grammarlens}

\subsection*{Your turn}
\begin{itemize}
\raggedright
  \item \emph{Say it:} \emph{pōrāṭuka}
  \item \emph{Say it:} \emph{pōrāṭuka}, once more
  \item \emph{Say it:} \emph{nāṭṭukār śatrukkaḷōṭŭ pōrāṭi.} — The local people fought the enemies.
  \item \emph{Recall:} say the Malayalam for a revolution, then the Malayalam for an invasion, then say \emph{pōrāṭuka} again
\end{itemize}

\begin{tcolorbox}[breakable,colback=teal!4,colframe=teal!35!black,title={Before you move on}]
What does \ml{പോരാടുക} mean? (\textquotedblleft{}To fight; to struggle\textquotedblright{}.) Say it once more.
\end{tcolorbox}
\section[\texorpdfstring{\ml{ചെറുത്തുനിൽക്കുക} (ceṟuttunilkkuka)}{ചെറുത്തുനിൽക്കുക (ceṟuttunilkkuka)}]{\ml{ചെറുത്തുനിൽക്കുക} (ceṟuttunilkkuka) — to resist; to hold out}

\label{lesson:ML-C401-ceruttunilkkuka}



\begin{quote}
Before the new one: say the Malayalam for an invasion, then the Malayalam for to fight.
\end{quote}

\subsection*{You'll want to know: \ml{ചെറുത്തുനിൽക്കുക}}
\textbf{\ml{ചെറുത്തുനിൽക്കുക}} — \emph{ceṟuttunilkkuka} — \textquotedblleft{}to resist; to hold out\textquotedblright{}.

\begin{quote}
\ml{ചെറിയ} \ml{ഗ്രാമം} \ml{ശത്രുസൈന്യത്തെ} \ml{ചെറുത്തുനിന്നു}.
\end{quote}

\begin{quote}
\emph{ceṟiya grāmaṁ śatrusainyatte ceṟuttuninnu.}
\end{quote}

\begin{quote}
The small village held out against the enemy army.
\end{quote}

Past: \ml{ചെറുത്തുനിന്നു}, \emph{ceṟuttuninnu}.

\ml{ചെറുത്ത്}, opposing, + \ml{നിൽക്കുക}, to stand: to stand and oppose.

\begin{cousinweb}
\textbf{\ml{ചെറുത്തുനിൽക്കുക}} (\emph{ceṟuttunilkkuka}) is a native Dravidian compound of \ml{ചെറുക്കുക}, \textquotedblleft{}to dam up, to oppose, to prevent\textquotedblright{} (DEDR 1980), and \ml{നിൽക്കുക} (\emph{nilkkuka}), \textquotedblleft{}to stand\textquotedblright{} (DEDR 3675).
\end{cousinweb}

\begin{grammarlens}[title={What you've built}]
One more word for telling what happened.
\end{grammarlens}

\subsection*{Your turn}
\begin{itemize}
\raggedright
  \item \emph{Say it:} \emph{ceṟuttunilkkuka}
  \item \emph{Say it:} \emph{ceṟuttunilkkuka}, once more
  \item \emph{Say it:} \emph{ceṟiya grāmaṁ śatrusainyatte ceṟuttuninnu.} — The small village held out against the enemy army.
  \item \emph{Recall:} say the Malayalam for an invasion, then the Malayalam for to fight, then say \emph{ceṟuttunilkkuka} again
\end{itemize}

\begin{tcolorbox}[breakable,colback=teal!4,colframe=teal!35!black,title={Before you move on}]
What does \ml{ചെറുത്തുനിൽക്കുക} mean? (\textquotedblleft{}To resist; to hold out\textquotedblright{}.) Say it once more.
\end{tcolorbox}
